% Pursuing Stacks (cotes 134-2 to 134-8): archive cote 134-7, site page 83 (PDF page 84), typescript page 525.
% typescript. Draft compared: ps6.tex:4154-4237.
% AI-transcribed by Claude Opus 5.5 (claude-opus-5-5, adaptive thinking, effort high), Message Batches API, batch ps4 (msgbatch_012Aoc1NuskATPsyBeKSYK7J), 2026-11-12; not yet checked by eye.
\scanpage{134-7/83}
\mspage{525}

(pointed) ``\deleted{abelian}\inserted{linear}ization'' of the schematic homotopy type X, \xout{whereas}
\xout{and to view the adjunction morphism} which maps into \deleted{it's abelianization$_X$}\inserted{the former}
\xout{Dually,} by the first adjunction morphism in (6 a). Dually, we may call
$\overset{\hand{\sim}}{S}(L\hand{\pi}.(X))$ the (pointed) ``co-\deleted{abelian}\inserted{linear}ization'' of X, it maps into X by the
first adjunction morphism in (6 b).

The \xout{first member in}\inserted[typed]{source} $L\overset{\hand{\sim}}{H}.(\overset{\hand{\sim}}{K}(L.))$ in the second adjunction map (6 a)
may be viewed as Eilenberg-Maclane type \inserted[typed]{(total)} homology, corresponding to the
chain complex L. \hand{,} (reducing to the \xout{familiar} more familiar $L\overset{\hand{\sim}}{H}.(\hand{\pi},n)$
when $L. =\hand{\pi}[n]$\unsure{\hand{)}} (but in the schematic sense of course ! ), whereas
the target $L\hand{\pi}.(\overset{\hand{\sim}}{S}(L.))$ in the second adjunction map (6 b) may be viewed
as ``total homotopy'' of a sphere-type space, reducing \xout{(}in case $L. =\hand{\pi}[n]$\xout{)}
to the total homotopy of the standard n-sphere over k, $S(n,\hand{\pi})$. In case
$k=\underline{Z}$ , we expect of course \xout{this to} the homology groups of this chain
complex to be the ordinary homotopy groups of the usual sphere $S^n$ -
cf. \hand{G)} below.

B) \uline{Inversion formulae}: two functorial isomorphisms
\[
(9)\qquad \hand{\Bigg\{}\begin{array}{ll}
a) & L\hand{\pi}.(\overset{\hand{\sim}}{K}(L.)) \xrightarrow{\hand{\sim}} L. \quad , \\
b) & L\overset{\hand{\sim}}{H}.(\overset{\hand{\sim}}{S}(L.)) \xrightarrow{\hand{\sim}} L. \quad .
\end{array}
\]

Maybe we should have begun with these, as they kind of fix the mea-%
ning of the two functors $\overset{\hand{\sim}}{K},\overset{\hand{\sim}}{S}$ in terms of $L\overset{\hand{\sim}}{H}.$ , $L\hand{\pi}.$ , which may be
viewed as embodying respectively the two main sets of invariants of a\xout{n}
homotopy type, namely homology and homotopy.

C) \uline{Hurewicz map} :
\[
(10)\qquad L\hand{\pi}.(X) \longrightarrow L\overset{\hand{\sim}}{H}.(X) \quad .
\]

We get such a map by applying $L\hand{\pi}.$ to the \deleted{abelian}\inserted{lineari}zation map in
(6 a), and using (9 a); symetrically, we may apply $L\overset{\hand{\sim}}{H}.$ to the
\xout{abelianization}\inserted[typed]{colineari}zation map (6 b), and use (9 b). We get a priori two maps
(10), and the statement is that these two maps are the same. Moreover
we have the all-important \uline{Hurewicz} \uline{theorem}: The first non-vanishing
homology object\hand{s} for $L\hand{\pi}.(X)$ and $L\overset{\hand{\sim}}{H}.(X)$ occur in the same dimension, n
say, and (10) induces an isomorphism for $H_n$ , \hand{an epimorphism for $H_{n+1}$.}

D) \uline{Exactness properties}: they can be stated shortly by saying that
in each one of the two pairs of adjoint functors (5), the left adjoint
one \xout{\unsure{commutes with cof}} respects cofibration sequences, whereas the right
adjoint respects fibration\xout{s} sequences. This may be detailed as four
exactness statements, one for each one of the four basic functors.

\margin{83}
